\section{Common Alfred Interface-Node Architecture}
\label{sec:common-platform}

Alfred does not develop one E2B platform and one unrelated Gateway platform. It develops \textbf{one common interface-node platform} that is configured, populated, or extended differently for the two customer architectures in Section~\ref{sec:exec-summary}.

\begin{diagram}
                       COMMON ALFRED NODE

              +---------------------------+
              | Protected 12/24 V Power   |
              | Voltage / Current Monitor |
              | MCU (SPI/UART/ADC/timers) |
              | Debug / Programming (SWD) |
              | Timing / Diagnostics      |
              | LAN8650/LAN8651 MAC-PHY   |
              | 10BASE-T1S                |
              +-------------+-------------+
                            |
               +------------+------------+
               |                         |
               v                         v

       PERIPHERAL INTERFACE        LEGACY-BUS INTERFACE
       (populated for E2B)         (populated for Gateway)

       GPIO   PWM   ADC            CAN   CAN-FD   LIN
       UART   SPI                  TX enable / bus monitor
       sensor / actuator I/O

               |                         |
               v                         v

       GREENFIELD E2B          BROWNFIELD GATEWAY
\end{diagram}

\subsection{What is genuinely shared}

\begin{center}
\small
\begin{tabularx}{\linewidth}{@{}p{3.4cm} X@{}}
\toprule
\textbf{Shared circuit block} & \textbf{Why it is identical on both variants} \\
\midrule
MCU and clocking & Same part, same clock tree, same firmware base image and bootloader across E2B and Gateway builds -- only the compiled peripheral driver set differs. \\
LAN8650/LAN8651 + magnetics/coupling & Same network silicon, same SPI wiring to the MCU, same termination-jumper layout, same PLCA node-ID strap scheme. This is the single most important shared block: it is what makes ``one platform, two roles'' true at the silicon level, not just at the marketing level. \\
Power input, protection, regulation & Same reverse-polarity, overcurrent, and transient-suppression front end; same 5\,V/3.3\,V regulator selection; same voltage/current monitoring circuit and MCU ADC channel assignment. \\
Reset, watchdog & Same supervisory reset IC, same watchdog strategy (hardware watchdog independent of MCU firmware). \\
Debug/programming & Same SWD header, pinout, and connector part number on every variant. \\
Timestamping / diagnostics & Same MCU timer/capture resource used for local event timestamping and the same UART/LED status scheme. \\
Manufacturing/test approach & Same bring-up test fixture and same first-boot provisioning procedure regardless of which interface section is populated. \\
PCB design rules & Same stackup, same impedance-controlled T1S trace rules, same connector keep-outs, same panelization strategy. \\
Expansion architecture & Same expansion header pinout and voltage rails, whether or not a given build populates it. \\
\bottomrule
\end{tabularx}
\end{center}

\subsection{What differs only by population/configuration}

The peripheral-interface section (GPIO/PWM/ADC/UART/SPI/sensor and actuator conditioning) and the legacy-bus section (CAN/CAN-FD/LIN transceivers, TX-enable gating, bus-state monitoring) occupy \emph{footprint-reserved but independently populated} real estate on the same PCB. An E2B build stuffs the peripheral section and leaves the legacy-bus footprints empty; a Gateway build does the reverse. Firmware follows the same split: one shared base image (network stack, manifest/health reporting, watchdog service) plus one of two thin role modules (peripheral I/O driver set vs. bus-translation driver set).

\subsection{Implementation strategy}

\begin{center}
\small
\begin{tabularx}{\linewidth}{@{}p{3.2cm} X@{}}
\toprule
\textbf{Option} & \textbf{Assessment} \\
\midrule
1. One common PCB, selective population & Lowest part-count and layout risk; one Gerber set, one stackup, one fab order feeds both roles. Cost: some wasted footprint area per build (acceptable at prototype volumes) and the legacy-bus and peripheral sections must both be laid out without conflicting for the same MCU pins/timers, which constrains pin assignment but is solvable at this pin count. \\
2. Common core board + daughtercards & Cleanest logical separation and best for a customer who wants to swap roles in the field, but adds a board-to-board connector, a second PCB program, and a second set of layout/EMI risks (connector impedance discontinuities matter on a T1S trace) -- more engineering surface area than this schedule can absorb before November. \\
3. Common reference design + closely related derivatives & Same core schematic/layout blocks reused across two \emph{physically separate but derivative} PCBs (copy-paste-modify rather than shared Gerbers). Avoids the population-conflict risk of Option 1 and the connector risk of Option 2, at the cost of maintaining two board files in parallel. \\
\bottomrule
\end{tabularx}
\end{center}

\begin{decision}[title={Recommendation}]
Use \textbf{Option 1 (one common PCB, selective population) for Rev~A and Rev~B}, where iteration speed and debugability matter most and board count must stay minimal. Move to \textbf{Option 3 (closely related derivatives from a common reference design) for Rev~C}, once the pin/footprint conflicts discovered in Rev~B are known well enough to split the E2B and Gateway variants into two clean, demo-quality boards without carrying unused footprint area into customer-facing hardware. At no point does the program build ``one E2B platform'' and ``one unrelated Gateway platform''; it builds one core architecture that is realized as Rev~A (learning board) $\rightarrow$ Rev~B (common node, both roles on one populated-selectively PCB) $\rightarrow$ Rev~C (E2B and Gateway derivatives of the same reference design).
\end{decision}
